\section{The Universal Mandate: Deep Evolutionary Origins of Sleep}\label{sec:universal_mandate}

Sleep is not a uniquely human or even mammalian trait. It appears to be a biological imperative for any nervous system. Sleep is ``nearly ubiquitous throughout the animal kingdom'' \citep{lesku2016evolutionary}. It has been characterized in insects like \textit{Drosophila} \citep{tian2023drosophila}, nematode worms, and even in jellyfish \citep{nathan2017jellyfish}, which possess a simple nerve net but no centralized brain. This evidence is bolstered by observations of sleep-like states in \textit{Hydra} \citep{kanaya2020hydra} and even sponges, which have no nervous system at all, suggesting they meet behavioral criteria for sleep \citep{bsj2020jellyfish}.

This suggests that sleep is a foundational, conserved process that evolved at the very advent of nervous systems. If these early-diverging lineages possess a sleep state, the common ancestor of all animals (the urmetazoan) likely had a primitive form of sleep \citep{lesku2016evolutionary}. Some biologists speculate sleep may even predate complex brains, with active wakefulness evolving later.

The most profound evidence for this universality comes from comparative biology, specifically the case of cephalopods. Octopuses (molluscs) and humans (vertebrates) sit on evolutionary branches that diverged over 550 million years ago \citep{medeiros2023octopus}. Yet, they have independently evolved a strikingly similar two-stage sleep architecture: a ``quiet sleep'' and an ``active sleep'' \citep{ramirezplascencia2023nature, medeiros2023octopus}. This ``active sleep'' in octopuses is a stunning parallel to human REM sleep. It involves body and eye twitching \citep{medeiros2023octopus, ramirezplascencia2023nature} and, most remarkably, the rapid, chaotic display of skin patterns (camouflage, hunting, threat) that reappear from their waking life \citep{ramirezplascencia2023nature}. While researchers cannot prove octopuses dream, the data is ``consistent with the idea'' \citep{medeiros2023octopus}.

If sleep-like states appear across virtually all animal phyla, then sleep is near-universal \citep{nathan2017jellyfish, SheinIdelson2016_Science_ReptileREM}.

This comparative evidence raises a profound question: If sleep is universal, is ``dreaming'' (i.e., offline generative neural simulation) also universal? The data suggests that animals from birds (who replay ``song'' in sleep) to bats (who replay ``flight paths'') experience forms of offline processing directly related to their waking life, even if we cannot access the phenomenology \citep{wilson1994reactivation, omer2025batreplay, suga2003, tian2023drosophila, lesku2016evolutionary}.

Sleep supports memory consolidation \citep{memory_function2010, raschborn2013}, regulates affective reactivity \citep{yoo2007}, and promotes abstraction/generalization \citep{stickgold2013, hidden_regularities2019}. Computationally, synaptic renormalization provides an evolutionary rationale \citep{price_of_plasticity2014}.
